\chapter{Marco Teórico}
\label{ch:marco-teorico}

\section{El Problema Multielectrónico y la Aproximación del Campo Medio}

El objetivo fundamental de la física atómica computacional es resolver la ecuación de Schrödinger independiente del tiempo para un sistema de \( N \) electrones interactuando con un núcleo fijo de carga \( Z \). El Hamiltoniano no relativista en unidades atómicas está dado por:

\begin{equation}
  \label{eq:hamiltonian-mb}
  \hat{H} = \sum_{i=1}^N \qty(-\frac{1}{2}\laplacian_i - \frac{Z}{r_i}) + \sum_{i<j}^N \frac{1}{\abs{\bm{r}_i-\bm{r}_j}}.
\end{equation}

Encontrar la función de onda exacta \( \Psi(\bm{x}_1, \bm{x}_2, \dots, \bm{x}_N) \) en este inmenso espacio de Hilbert de \( 3N \) dimensiones espaciales es analítica y numéricamente intratable para \( N \geq 3 \). El obstáculo principal radica en el término de repulsión interelectrónica (el operador de dos cuerpos), el cual acopla todas las coordenadas de las partículas, impidiendo la separación de variables. 

Para sortear esta barrera, introducimos la aproximación de campo medio o Hartree-Fock. En esta aproximación, se asume que cada electrón se mueve independientemente de los demás, sintiendo únicamente el potencial electrostático esférico promedio generado por el núcleo y la nube de carga del resto de los electrones. 

\section{El Determinante de Slater y la Antisimetría}

Debido a que los electrones son fermiones, la mecánica cuántica exige que la función de onda total del sistema sea estrictamente antisimétrica bajo el intercambio de cualquier par de partículas (Principio de exclusión de Pauli). Matemáticamente, la forma más sencilla de construir un estado antisimétrico a partir de orbitales de partículas individuales (espín-orbitales) es mediante un determinante, conocido como el \textbf{Determinante de Slater}:

\begin{equation}
  \Psi(\bm{x}_1, \bm{x}_2, \dots, \bm{x}_N) = \frac{1}{\sqrt{N!}}
  \vmqty{
    \phi_1(\bm{x}_1) & \phi_2(\bm{x}_1) & \cdots & \phi_N(\bm{x}_1) \\
    \phi_1(\bm{x}_2) & \phi_2(\bm{x}_2) & \cdots & \phi_N(\bm{x}_2) \\
    \vdots & \vdots & \ddots & \vdots \\
    \phi_1(\bm{x}_N) & \phi_2(\bm{x}_N) & \cdots & \phi_N(\bm{x}_N)
  }.
\end{equation}

El objetivo central del método de Hartree-Fock es evaluar el funcional de energía del sistema asumiendo que el estado es un único determinante de Slater:

\begin{equation}
  \label{eq:energy-func}
  E = \ev{\hat{H}}{\Psi},
\end{equation}

y encontrar variacionalmente el conjunto de espín-orbitales \( \{ \phi_i \} \) que minimicen esta energía, manteniendo la condición de ortonormalidad \( \braket{\phi_i}{\phi_j} = \delta_{ij} \).

\section{Suficiencia del Modelo de Dos Electrones}

Al expandir el valor esperado de la energía \eqref{eq:energy-func} sobre el determinante de Slater de $N$ cuerpos, el problema parece, a primera vista, colosalmente complejo por sus expansiones factoriales. Sin embargo, observemos con detenimiento la estructura del Hamiltoniano atómico \eqref{eq:hamiltonian-mb}: está compuesto exclusivamente por operadores de \textbf{un cuerpo} (energía cinética y atracción nuclear) y operadores de \textbf{dos cuerpos} (repulsión coulómbica entre pares de electrones). 

No existen operadores fundamentales de tres o más cuerpos en el Hamiltoniano electrostático. Esta profunda simplificación física dictamina que, sin importar si estamos resolviendo el litio ($Z=3$) o el oganesón ($Z=118$), el acoplamiento directo siempre ocurre exclusivamente entre pares de orbitales. Por lo tanto, aislar el comportamiento de un sistema de solo dos electrones es suficiente para derivar las reglas de interacción generales, las cuales posteriormente escalaremos como una sumatoria por pares para todo el átomo.

Consideremos entonces un sistema abstracto de dos electrones descritos por los espín-orbitales \( \ket{a} \) y \( \ket{b} \). El determinante se reduce a una matriz de $2 \times 2$:

\begin{equation}
  \ket{\Psi_{ab}} = \frac{1}{\sqrt{2}} [\phi_a(\bm{x}_1)\phi_b(\bm{x}_2) - \phi_b(\bm{x}_1)\phi_a(\bm{x}_2)].
\end{equation}

Al evaluar el operador perturbativo de repulsión \( \frac{1}{r_{12}} \), obtenemos:

\begin{equation}
  \ev{\frac{1}{r_{12}}}{\Psi_{ab}} = \frac{1}{2} \int \dd{\bm{x}_1} \dd{\bm{x}_2} \abs{\phi_a(\bm{x}_1)\phi_b(\bm{x}_2) - \phi_b(\bm{x}_1) \phi_a(\bm{x}_2) }^2 \frac{1}{r_{12}}.
\end{equation}

Al desarrollar el binomio al cuadrado, se revelan dos tipos cualitativamente distintos de interacciones electrónicas, producto directo de la imposición matemática de la antisimetría:

\begin{itemize}
\item \textbf{Interacción Directa (Integrales de Coulomb, \(J_{ab}\)):}
  Provienen de los términos cuadrados perfectos. Reflejan la electrostática clásica: la repulsión de Coulomb entre la nube de densidad de carga del electrón $a$ y la nube de densidad del electrón $b$. 
  \begin{equation}
  J_{ab} = \int \dd{\bm{r}_1} \dd{\bm{r}_2} \abs{\psi_a(\bm{r}_1)}^2 \frac{1}{r_{12}} \abs{\psi_b(\bm{r}_2)}^2.
  \end{equation}

\item \textbf{Interacción de Interferencia (Integrales de Intercambio, \(K_{ab}\)):}
  Surgen exclusivamente del término cruzado del binomio. Es un efecto puramente cuántico sin equivalente clásico, generado por el solapamiento indistinguible de los estados. Este término aparece únicamente si ambos electrones poseen el mismo estado de espín, creando un "Hueco de Fermi": una repulsión topológica que mantiene alejados a los electrones con espines paralelos más allá de su repulsión de carga.
  \begin{equation}
  K_{ab} = \delta_{\sigma_a \sigma_b} \int \dd{\bm{r}_1} \dd{\bm{r}_2} \psi_a^*(\bm{r}_1)\psi_b^*(\bm{r}_2)\frac{1}{r_{12}}\psi_b(\bm{r}_1)\psi_a(\bm{r}_2).
  \end{equation}
\end{itemize}

Gracias a las Reglas de Slater-Condon, se demuestra formalmente que la energía total de repulsión para el determinante de $N$ cuerpos es simplemente la suma de las interacciones acopladas de los pares:
\begin{equation}
  E_{ee} = \sum_{a=1}^N \sum_{b>a}^N \qty( J_{ab} - K_{ab} ).
\end{equation}

\section{Expansión Multipolar y la Necesidad Iterativa}

En un campo con simetría central (donde la base orbital $\phi_{a}$ se puede separar en una parte radial $P_a(r)$ y otra angular asociada a armónicos esféricos), el operador $\frac{1}{r_{12}}$ se puede descomponer mediante armónicos esféricos. Esto nos permite separar analíticamente la dependencia angular (geométrica) de la física radial. (Para detalles exhaustivos sobre tensores irreducibles, consúltese el \textbf{Apéndice \ref{ap:tensores}}). 

La repulsión interelectrónica se colapsa entonces en lo que definimos como \textbf{Integrales Radiales de Slater}, denotadas como $R^k(abcd)$. Estas integrales miden el acoplamiento radial real entre orbitales mediante potenciales esféricos efectivos de orden multipolar $k$:

\begin{equation}
  \label{eq:slater_k}
  R^k(abcd) = \int_0^\infty P_a(r_1) \qty[ \int_0^\infty \frac{r_{<}^k}{r_{>}^{k+1}} P_b(r_2) P_d(r_2) \dd{r_2} ] P_c(r_1) \dd{r_1}.
\end{equation}

Estas integrales (que engloban tanto el acoplamiento Coulombiano directo como el Intercambio) dictan la estructura fina de los niveles de energía. Sin embargo, analizar la Ecuación \eqref{eq:slater_k} revela el monumental obstáculo del límite de Hartree-Fock.

Para determinar las funciones radiales $P(r)$ exactas que minimicen la energía, debemos someterlas a la acción de un "Operador de Fock" efectivo que contiene al potencial electrostático de los demás electrones. Pero, como se observa en la formulación, el potencial en sí mismo está construido a partir de los orbitales radiales $P_b$ y $P_d$ incógnitos. \textbf{Las funciones de onda definen al potencial, y el potencial define a las funciones de onda.}

Esta severa no-linealidad e interdependencia condena cualquier esfuerzo de solución analítica directa. Nos sitúa ante la inevitable necesidad metodológica de recurrir a ciclos de convergencia numérica, donde la única forma de descifrar la estructura atómica es forzando iterativamente al sistema hasta que su estado de respuesta se equipare a su estado de entrada. Esta paradoja de autogeneración establece la premisa y motivación central para transitar al siguiente capítulo: el desarrollo de un algoritmo de Campo Autoconsistente (SCF).
